% GENERATED FILE. Edit canonical lessons, then run npm run generate:books.
% canonical-source-hash: fnv1a64:5555c1c372d54430
% canonical-lessons: TE-C452-revu, TE-C452-nauka, TE-C452-navikudu, TE-C452-langaru, TE-C452-dipastambham

\chapter{Ships and Harbours}
\label{ch:te-ships-and-harbours}
\hlchaptermodality{452}

\begin{chapteropening}
\textbf{By the end of this chapter:} \emph{I can say a harbour, a ship, a sailor, an anchor and a lighthouse.}

Complete the last lesson of chapter 452: I can say a harbour, a ship, a sailor, an anchor and a lighthouse.
\end{chapteropening}

\section[\texorpdfstring{\te{రేవు} (rēvu)}{రేవు (rēvu)}]{\te{రేవు} (rēvu) — a harbour; a landing place}

\label{lesson:TE-C452-revu}



\begin{quote}
Before the new one: say the Telugu for going on foot, then the Telugu for a wayfarer.
\end{quote}

\subsection*{You'll want to know: \te{రేవు}}
\textbf{\te{రేవు}} — \emph{rēvu} — \textquotedblleft{}a harbour; a landing place\textquotedblright{}.

\begin{quote}
\te{సాయంత్రం} \te{పడవలు} \te{రేవుకు} \te{తిరిగి} \te{వచ్చాయి}.
\end{quote}

\begin{quote}
\emph{sāyantraṁ paḍavalu rēvuku tirigi vaccāyi.}
\end{quote}

\begin{quote}
In the evening the boats came back to the harbour.
\end{quote}

A harbour for ships, and also the steps down to a river where people bathe and boats land.

\begin{cousinweb}
\textbf{\te{రేవు}} (\emph{rēvu}) is native Dravidian: DEDR 516 lists Telugu \emph{rēvu}, \textquotedblleft{}a landing place, a port, a ford\textquotedblright{}, in the family of Tamil \emph{iṟaṅku}, \textquotedblleft{}to descend, to disembark\textquotedblright{}.
\end{cousinweb}

\begin{grammarlens}[title={What you've built}]
One more word for getting around.
\end{grammarlens}

\subsection*{Your turn}
\begin{itemize}
\raggedright
  \item \emph{Say it:} \emph{rēvu}
  \item \emph{Say it:} \emph{rēvu}, once more
  \item \emph{Say it:} \emph{sāyantraṁ paḍavalu rēvuku tirigi vaccāyi.} — In the evening the boats came back to the harbour.
  \item \emph{Recall:} say the Telugu for going on foot, then the Telugu for a wayfarer, then say \emph{rēvu} again
\end{itemize}

\begin{tcolorbox}[breakable,colback=teal!4,colframe=teal!35!black,title={Before you move on}]
What does \te{రేవు} mean? (\textquotedblleft{}A harbour; a landing place\textquotedblright{}.) Say it once more.
\end{tcolorbox}
\section[\texorpdfstring{\te{నౌక} (nauka)}{నౌక (nauka)}]{\te{నౌక} (nauka) — a ship; a vessel}

\label{lesson:TE-C452-nauka}



\begin{quote}
Before the new one: say the Telugu for a wayfarer, then the Telugu for a harbour.
\end{quote}

\subsection*{You'll want to know: \te{నౌక}}
\textbf{\te{నౌక}} — \emph{nauka} — \textquotedblleft{}a ship; a vessel\textquotedblright{}.

\begin{quote}
\te{పెద్ద} \te{నౌక} \te{సముద్రంలో} \te{దూరంగా} \te{కనిపించింది}.
\end{quote}

\begin{quote}
\emph{pedda nauka samudraṁlō dūraṅgā kanipiñcindi.}
\end{quote}

\begin{quote}
A big ship appeared far out at sea.
\end{quote}

Plural: \te{నౌకలు}, \emph{naukalu}.

A formal word. Plural \te{నౌకలు}.

\begin{cousinweb}
\textbf{\te{నౌక}} (\emph{nauka}) is borrowed from Sanskrit \emph{naukā}, \textquotedblleft{}a boat, a ship\textquotedblright{} (Monier-Williams).
\end{cousinweb}

\begin{grammarlens}[title={What you've built}]
One more word for getting around.
\end{grammarlens}

\subsection*{Your turn}
\begin{itemize}
\raggedright
  \item \emph{Say it:} \emph{nauka}
  \item \emph{Say it:} \emph{nauka}, once more
  \item \emph{Say it:} \emph{pedda nauka samudraṁlō dūraṅgā kanipiñcindi.} — A big ship appeared far out at sea.
  \item \emph{Recall:} say the Telugu for a wayfarer, then the Telugu for a harbour, then say \emph{nauka} again
\end{itemize}

\begin{tcolorbox}[breakable,colback=teal!4,colframe=teal!35!black,title={Before you move on}]
What does \te{నౌక} mean? (\textquotedblleft{}A ship; a vessel\textquotedblright{}.) Say it once more.
\end{tcolorbox}
\section[\texorpdfstring{\te{నావికుడు} (nāvikuḍu)}{నావికుడు (nāvikuḍu)}]{\te{నావికుడు} (nāvikuḍu) — a sailor}

\label{lesson:TE-C452-navikudu}



\begin{quote}
Before the new one: say the Telugu for a harbour, then the Telugu for a ship.
\end{quote}

\subsection*{You'll want to know: \te{నావికుడు}}
\textbf{\te{నావికుడు}} — \emph{nāvikuḍu} — \textquotedblleft{}a sailor\textquotedblright{}.

\begin{quote}
\te{నా} \te{మామయ్య} \te{నావికుడిగా} \te{పనిచేశాడు}.
\end{quote}

\begin{quote}
\emph{nā māmayya nāvikuḍigā panicēśāḍu.}
\end{quote}

\begin{quote}
My uncle worked as a sailor.
\end{quote}

\te{నావ}, a boat, + -\te{ఇకుడు}, one who works with it.

\begin{cousinweb}
\textbf{\te{నావికుడు}} (\emph{nāvikuḍu}) is borrowed from Sanskrit \emph{nāvika}, \textquotedblleft{}a sailor, a helmsman\textquotedblright{} (CDIAL 7082), from \emph{nāvā}, \textquotedblleft{}a boat\textquotedblright{}.
\end{cousinweb}

\begin{grammarlens}[title={What you've built}]
One more word for getting around.
\end{grammarlens}

\subsection*{Your turn}
\begin{itemize}
\raggedright
  \item \emph{Say it:} \emph{nāvikuḍu}
  \item \emph{Say it:} \emph{nāvikuḍu}, once more
  \item \emph{Say it:} \emph{nā māmayya nāvikuḍigā panicēśāḍu.} — My uncle worked as a sailor.
  \item \emph{Recall:} say the Telugu for a harbour, then the Telugu for a ship, then say \emph{nāvikuḍu} again
\end{itemize}

\begin{tcolorbox}[breakable,colback=teal!4,colframe=teal!35!black,title={Before you move on}]
What does \te{నావికుడు} mean? (\textquotedblleft{}A sailor\textquotedblright{}.) Say it once more.
\end{tcolorbox}
\section[\texorpdfstring{\te{లంగరు} (laṅgaru)}{లంగరు (laṅgaru)}]{\te{లంగరు} (laṅgaru) — an anchor}

\label{lesson:TE-C452-langaru}



\begin{quote}
Before the new one: say the Telugu for a ship, then the Telugu for a sailor.
\end{quote}

\subsection*{You'll want to know: \te{లంగరు}}
\textbf{\te{లంగరు}} — \emph{laṅgaru} — \textquotedblleft{}an anchor\textquotedblright{}.

\begin{quote}
\te{నావికులు} \te{రేవులో} \te{లంగరు} \te{వేశారు}.
\end{quote}

\begin{quote}
\emph{nāvikulu rēvulō laṅgaru vēśāru.}
\end{quote}

\begin{quote}
The sailors dropped anchor in the harbour.
\end{quote}

\te{లంగరు} \te{వేయు}, to drop anchor; \te{లంగరు} \te{ఎత్తు}, to weigh anchor.

\begin{cousinweb}
\textbf{\te{లంగరు}} (\emph{laṅgaru}) is usually traced to Persian \emph{langar}, \textquotedblleft{}an anchor\textquotedblright{}, which reached the coasts of India with sea trade.
\end{cousinweb}

\begin{grammarlens}[title={What you've built}]
One more word for getting around.
\end{grammarlens}

\subsection*{Your turn}
\begin{itemize}
\raggedright
  \item \emph{Say it:} \emph{laṅgaru}
  \item \emph{Say it:} \emph{laṅgaru}, once more
  \item \emph{Say it:} \emph{nāvikulu rēvulō laṅgaru vēśāru.} — The sailors dropped anchor in the harbour.
  \item \emph{Recall:} say the Telugu for a ship, then the Telugu for a sailor, then say \emph{laṅgaru} again
\end{itemize}

\begin{tcolorbox}[breakable,colback=teal!4,colframe=teal!35!black,title={Before you move on}]
What does \te{లంగరు} mean? (\textquotedblleft{}An anchor\textquotedblright{}.) Say it once more.
\end{tcolorbox}
\section[\texorpdfstring{\te{దీపస్తంభం} (dīpastambhaṁ)}{దీపస్తంభం (dīpastambhaṁ)}]{\te{దీపస్తంభం} (dīpastambhaṁ) — a lighthouse}

\label{lesson:TE-C452-dipastambham}



\begin{quote}
Before the new one: say the Telugu for a sailor, then the Telugu for an anchor.
\end{quote}

\subsection*{You'll want to know: \te{దీపస్తంభం}}
\textbf{\te{దీపస్తంభం}} — \emph{dīpastambhaṁ} — \textquotedblleft{}a lighthouse\textquotedblright{}.

\begin{quote}
\te{రాత్రి} \te{దీపస్తంభం} \te{వెలుగు} \te{నౌకలకు} \te{దారి} \te{చూపిస్తుంది}.
\end{quote}

\begin{quote}
\emph{rātri dīpastambhaṁ velugu naukalaku dāri cūpistundi.}
\end{quote}

\begin{quote}
At night the light of the lighthouse shows the ships the way.
\end{quote}

\te{దీపం}, a lamp, + \te{స్తంభం}, a pillar: a pillar with a lamp.

\begin{cousinweb}
\textbf{\te{దీపస్తంభం}} (\emph{dīpastambhaṁ}) is a compound of two words borrowed from Sanskrit: \emph{dīpa}, \textquotedblleft{}a lamp\textquotedblright{} (CDIAL 6348), and \emph{stambha}, \textquotedblleft{}a pillar, a post\textquotedblright{} (CDIAL 13682).
\end{cousinweb}

\begin{grammarlens}[title={What you've built}]
One more word for getting around.
\end{grammarlens}

\subsection*{Your turn}
\begin{itemize}
\raggedright
  \item \emph{Say it:} \emph{dīpastambhaṁ}
  \item \emph{Say it:} \emph{dīpastambhaṁ}, once more
  \item \emph{Say it:} \emph{rātri dīpastambhaṁ velugu naukalaku dāri cūpistundi.} — At night the light of the lighthouse shows the ships the way.
  \item \emph{Recall:} say the Telugu for a sailor, then the Telugu for an anchor, then say \emph{dīpastambhaṁ} again
\end{itemize}

\begin{tcolorbox}[breakable,colback=teal!4,colframe=teal!35!black,title={Before you move on}]
What does \te{దీపస్తంభం} mean? (\textquotedblleft{}A lighthouse\textquotedblright{}.) Say it once more.
\end{tcolorbox}
